% Results Section — v2 (2026-04-17)
% Target: eLife / Current Biology / Nature Communications
% Figure assets: Figures/generated/ (imported from ICML_workshop/icml2026/paperbanana_inputs)
%   fig1a_output.png        — Pipeline schematic (Fig 1a)
%   fig1_panels_bcd.pdf     — LOCO vulnerability profiles (Fig 1b-d)
%   fig2_output.png         — Correction feasibility / arc collapse (Fig 2)
%   fig_hc_specificity.pdf  — HC FPR calibration (Fig S2; supplement)
%
% Structure (5 subsections):
%   R1. Behavioral phenotype and CVD severity
%   R2. Color discrimination is preserved but continuous interpolation is impaired (Q1,Q2)
%   R3. Low-dimensional distortion fields recovered per CVD individual (Q3 part 1)
%   R4. Convergent detection, divergent correction (Q3 part 2 — filter substrate)
%   R5. Specificity calibration and cross-criterion validation (methodological honesty)

\section{Results}
\label{sec:results}

\subsection{Behavioral phenotype and CVD severity}
\label{sec:results:behavior}

Three CVD participants were recruited alongside seven healthy controls
(HC). All CVD participants were male and completed the 14-plate Ishihara
pseudoisochromatic test~\citep{ishihara1917}. Sub-08 was classified as
moderate deuteranomalous (8/14 plates misread, protan--deutan
discrimination plate indicating deutan axis), Sub-09 as moderate--severe
protanomalous (11/14 plates misread, protan axis), and Sub-10 as
near-normal deutan (2/14 plates misread; included as a within-study
specificity control rather than as a symptomatic case). Same-day
8-alternative forced-choice identification of the eight scanner stimuli
confirmed the Ishihara categorisation: HC accuracy was
$0.94 \pm 0.04$ (mean $\pm$ SD), Sub-10 $= 0.88$, Sub-09 $= 0.62$,
Sub-08 $= 0.71$ (chance $= 0.125$). Errors in Sub-08 clustered on the
red--green confusion axis (colours c2/c3 vs.\ c7; mean confusion 28\%);
errors in Sub-09 clustered along the full L--M opponent range with
additional misidentification of magenta (c8).

Together with the genetic range of L/M cone opsin polymorphisms
(2--12~nm spectral separation;~\citealp{neitz2011, deeb2005}) and
population estimates of 6--8\% deutan-type and 1--2\% protan-type
anomalous trichromacy in males~\citep{bosten2019}, these three
participants span the clinically relevant severity spectrum from
near-normal to moderate--severe, with one dichromat-adjacent protan case
(Sub-09) and one mild deutan case (Sub-08). All subsequent neural
analyses interpret each CVD participant as an individual case, following
the single-case statistical framework of \citet{crawford1998}.

\subsection{Color discrimination is preserved but continuous
interpolation is impaired}
\label{sec:results:dissociation}

We first asked whether CVD alters the neural representation of color at
all (Q1), and if so whether it affects category-level discrimination or
continuous relational geometry (Q2). Two cross-validation schemes
applied to the hV4 forward encoding model dissociate these two
properties (\cref{fig:main}a): leave-one-run-out (LORO) probes whether
the eight isoluminant hues can be \emph{classified} from held-out runs,
while leave-one-colour-out (LOCO) probes whether the encoder
\emph{interpolates} to a hue never seen in training.

\paragraph{Discrimination (LORO) is group-equivalent.}
Classification accuracy in hV4 exceeded the 0.125 chance rate in every
individual (HC $= 0.68 \pm 0.11$, CVD $= 0.64 \pm 0.09$, mean $\pm$ SD
across 6 runs). A between-group permutation test on accuracy showed no
HC--CVD difference at any ROI (hV4 $p = 0.668$; V1 $p = 0.542$; V2
$p = 0.611$; 10{,}000 label permutations), and the same pattern held
when decoding was run on SRM-aligned features
(\cref{tab:loro,app:tables}). CVD therefore does not abolish the
category-level identity code for the eight scanner hues.

\paragraph{Interpolation (LOCO) is selectively impaired in CVD.}
The picture changes sharply under LOCO. HC participants achieve above-chance
hue interpolation in hV4 ($\rho = 0.42 \pm 0.14$ between held-out and
predicted responses; cf.~\citealp{brouwer2009}), whereas both CVD
participants fall near floor (Sub-08 $\rho = 0.08$; Sub-09 $\rho = 0.12$;
Hedges' $g_{\text{HC vs.\ CVD}} = 1.61$ in V1, $g = 1.85$ in V2, $g = 1.34$
in hV4; all $p_{\text{perm}} < 0.03$; \cref{fig:main}b--c, gray bars).
Sub-10, the near-normal control, shows HC-like interpolation
($\rho = 0.38$; \cref{fig:main}d). Critically, the \emph{per-colour}
vulnerability vector $\mathbf{v} \in \mathbb{R}^{8}$ is individually
structured: Sub-08's largest errors lie on the c2/c3 (orange/yellow)
and c7 (purple) channels, whereas Sub-09's errors cluster on c5/c6
(cyan/blue) with magenta (c8) as a further outlier. This per-colour
fingerprint, invisible to the group-level LORO test, is the geometry-sensitive
phenotype on which all subsequent modelling operates
(\cref{sec:results:inference}).

The LORO $\to$ LOCO dissociation is conceptually important: the same
hV4 code that supports category-level readout fails to support
neighbour-based inference in CVD. This is the precise regime in which
continuous downstream tasks---hue naming, traffic-light perception,
display rendering---are affected, and it identifies the representational
substrate on which an assistive filter must act.

\begin{figure*}[tb]
  \centering
  \includegraphics[width=\textwidth]{Figures/generated/fig1a_output.png}\\[0.6em]
  \includegraphics[width=0.95\textwidth]{Figures/generated/fig1_panels_bcd.pdf}
  \caption{\textbf{Inferring subject-specific distortion fields from
  structured neural representations.}
  \textbf{(a)} Pipeline. Procrustes-aligned fMRI amplitudes
  ($6~\text{runs} \times 8~\text{hues} \times V$ voxels in hV4) feed a
  forward encoder with $K{=}3$ half-wave-rectified squared sinusoidal
  channels. Leave-one-colour-out cross-validation yields the
  \emph{vulnerability profile} $\mathbf{v} \in \mathbb{R}^{8}$, which is
  inverted to $\leq$2 DOF distortion parameters (retinal
  Machado / retinal+cortical R+C / cortical 2-component). Pre-image
  inversion produces a per-hue correction filter whose bijectivity is
  assessed before behavioural validation.
  \textbf{(b--d)} Observed vs.\ best-model LOCO vulnerability per
  participant. Gray bars: empirical $\mathbf{v}_{\text{CVD}}$;
  coloured line: best-family simulated $\mathbf{v}_{\text{sim}}$.
  \textbf{(b)} Sub-08 (deutan): R+C model,
  $\Delta\lambda = 2.0$~nm, $g = +2.25$, $\rho = 0.857$, $p = 0.005$.
  \textbf{(c)} Sub-09 (protan): Machado model,
  $\Delta\lambda = 13.5$~nm, $\rho = 0.762$, $p = 0.018$.
  \textbf{(d)} Sub-10 (near-normal): no cone-level model reaches
  significance, $p = 0.559$ — specificity confirmed. $p$ values are
  $8! = 40{,}320$ exact permutations of colour labels on Spearman
  $\rho$.}
  \label{fig:main}
\end{figure*}

\subsection{Low-dimensional distortion fields recovered per CVD
individual}
\label{sec:results:inference}

Having established that CVD distorts continuous hue geometry
(\cref{sec:results:dissociation}), we next asked whether this
distortion is parsimoniously describable (Q3). For each participant we
fitted three mechanistic models by grid search, minimising the
composite loss of \cref{eq:loss}: the 1-DOF Machado cone shift
($\Delta\lambda$), the 2-DOF retinal+cortical R+C model
($\Delta\lambda, g$), and the 2-DOF cortical 2-component angular
dilation ($\beta_s, \beta_c$). Because R+C strictly contains Machado
as its $g{=}0$ limit, we applied a closed-testing procedure: the more
complex model was only reported if the simpler one was not
significant.

\paragraph{Sub-09 (moderate protanomaly).}
The 1-DOF Machado model with $\Delta\lambda = 13.5$~nm captured the
hV4 vulnerability profile at $\rho = 0.762$,
$p_{\text{perm}} = 0.018$ (\cref{fig:main}c). The shift magnitude is
within the established biophysical range for moderate--severe
protanomaly~\citep{machado2009, neitz2011}, and closed-testing
escalation to R+C yielded $g_{\ast} = 0.0$ at the optimum, i.e.\ no
cortical augmentation beyond Machado. A purely retinal description is
therefore sufficient for Sub-09.

\paragraph{Sub-08 (mild deuteranomaly).}
Here Machado alone was only trend-level ($\rho = 0.619$,
$p = 0.058$). Closed-testing escalated to R+C, which converged at
$\Delta\lambda = 2.0$~nm, $g = +2.25$, $\rho = 0.857$,
$p = 0.005$ (\cref{fig:main}b; \cref{tab:main}). The retinal shift is
small --- consistent with the mild Ishihara phenotype --- while the
cortical gain parameter is large and positive, suggesting that mild
peripheral anomaly in this participant is accompanied by a
non-negligible cortical re-weighting of the residual L--M opponent
signal. Because $g > 0$ points away from the "exact compensation"
regime ($g = -1$; \citealp{tregillus2021}) and into apparent
over-amplification, we treat it as an effective description rather
than a biophysical claim (\cref{sec:discussion}).

\paragraph{Sub-10 (near-normal, deutan-axis).}
No cone-level model reached significance
($\rho = -0.048$, $p = 0.559$; $\Delta\lambda \approx 0$ at the
optimum; \cref{fig:main}d, \cref{tab:main}). This within-study
specificity control confirms that the pipeline does not fabricate a
distortion where none is present at the neural level despite the mild
Ishihara signal.

\begin{table}[tb]
  \caption{Inferred distortion parameters in hV4 LOCO (primary ROI).
  Closed testing: the 1-DOF Machado model is tested first; the 2-DOF
  R+C or 2-component models are reported when Machado is not
  significant. $p$-values are $8! = 40{,}320$ exact permutations on
  the colour-label Spearman $\rho$.}
  \label{tab:main}
  \centering
  \small
  \begin{tabular}{llccccl}
    \toprule
    Subject & Winning model & DOF &
    $\Delta\lambda$~(nm) &
    Cortical & $\rho$ & $p_{\text{perm}}$ \\
    \midrule
    Sub-08 (deutan)   & R+C       & 2 & 2.0  & $g = +2.25$                         & 0.857 & \textbf{0.005} \\
    Sub-09 (protan)   & Machado   & 1 & 13.5 & ---                                 & 0.762 & \textbf{0.018} \\
    Sub-10 (control)  & (none)    & --- & $\sim$0 & ---                            & $-0.048$ & 0.559 \\
    \bottomrule
  \end{tabular}
\end{table}

\paragraph{Cortical-level parameters converge on a common S-cone
expansion angle.}
The complementary 2-component cortical model reaches significance for
both CVD participants (Sub-08: $\beta_s = 38^\circ$, $\beta_c = -14^\circ$,
$\rho = 0.881$, $p = 0.004$; Sub-09: $\beta_s = 6^\circ$,
$\beta_c = -22^\circ$, $\rho = 0.690$, $p = 0.035$;
\cref{tab:full,app:supporting}). When fitted not against LOCO but
against the pairwise $\Delta$RDM in V1, $\beta_s$ converges at
$\sim 21.5^\circ$ across both CVD participants, matching the
$21.4^\circ$ behavioural S-cone re-weighting reported by
\citet{emery2021} in an independent anomalous-trichromat cohort. This
cross-cohort, cross-modality convergence is a non-trivial validation:
the same S-cone expansion angle emerges from our fMRI V1 geometry and
from Emery et~al.'s hue-scaling behaviour, suggesting that anomalous
trichromats recruit a shared compensatory axis irrespective of subtype.

\subsection{Convergent detection, divergent correction: a filter
substrate}
\label{sec:results:filter}

\cref{sec:results:inference} establishes \emph{what} is distorted per
participant. Because the question Q3 concerns whether such distortion
can be corrected, we next inverted each fitted distortion field into a
stimulus-space pre-image filter and examined its behaviour.

\paragraph{Retinal and cortical models agree on \emph{who} is
distorted, but disagree on \emph{how} to correct.}
Every model family --- Machado, R+C, 2-component --- independently
flags Sub-08 and Sub-09 as distorted and Sub-10 as null
(\cref{tab:main,tab:full}). However, the prescribed corrections point
in opposite directions. For Sub-08, the R+C-derived per-hue
correction vector and the 2-component-derived correction vector have
cosine similarity $-0.18$ and sign agreement $3/8$. Both models
explain the neural profile comparably well
($\rho_{\text{R+C}} = 0.857$ vs.\ $\rho_{\text{2-comp}} = 0.881$), so
model fit alone cannot adjudicate between them. This divergence is the
central falsifiable prediction of the study: the two models make
different quantitative claims about which displayed hues require
which angular nudge, and those claims can be tested behaviourally.

\paragraph{Correction feasibility depends on distortion severity and
modelling level.}
The cone-shift model predicts an opponent-space arc compression that
scales with $\Delta\lambda$ (\cref{fig:feasibility}). For Sub-08's
small shift ($\Delta\lambda = 2.0$~nm) the arc remains wide and all 8
hues admit an exact pre-image under R+C (mean $|\delta| = 23.5^\circ$,
max $= 42.9^\circ$, residual $< 0.001^\circ$; \cref{fig:feasibility}a).
For Sub-09's moderate shift ($\Delta\lambda = 13.5$~nm) the opponent
arc compresses from $360^\circ$ to $\sim 96^\circ$: three hues
(c4 green, c5 cyan, c6 blue) converge onto a single perceived angle
($\sim 282^\circ$) and exact Machado-level inversion becomes
structurally impossible --- only 4/8 hues have an exact pre-image, and
the remaining 4 require a separation-optimised approximate solution
(min separation $1.03^\circ \to 5.76^\circ$, a $5.6\times$
improvement but still well below the healthy inter-hue spacing of
$\sim 70^\circ$; \cref{fig:feasibility}b, upper row).

\paragraph{The cortical-level model preserves bijectivity by
construction.}
Under the 2-component angular dilation model, distortion is a smooth
rotation in opponent hue space and its inverse is therefore guaranteed
to exist. Applying this inverse to Sub-09 recovers all 8 hues exactly
(mean $|\delta| = 20.1^\circ$, max $= 48.1^\circ$, residual
$< 0.001^\circ$; \cref{fig:feasibility}b, lower row). For Sub-08 the
cortical-model pre-image is likewise exact for 8/8 hues. The cortical
model is therefore the only one that admits an exact stimulus-space
correction for \emph{both} severity levels, and it is the only model
to pass dual-criterion validation (LOCO \emph{and} $\Delta$RDM; see
\cref{sec:results:specificity}). It is on this basis that we designate
the 2-component model as the \emph{filter substrate} for the behavioural
experiment in \cref{sec:results:behavioralfilter}, while retaining the
R+C retinal model as a rival hypothesis whose divergent correction
vector provides the falsifier.

\begin{figure*}[tb]
  \centering
  \includegraphics[width=0.95\textwidth]{Figures/generated/fig2_output.png}
  \caption{\textbf{Correction feasibility depends jointly on
  distortion severity and the modelling level.}
  \textbf{(a)} Sub-08 (mild deuteranomaly, R+C fit): the opponent hue
  arc remains wide, all 8 displayed hues undergo small angular
  shifts (hollow dots), and each admits an exact pre-image correction
  (gray arrows; $< 0.001^\circ$ residual; 8/8 exact).
  \textbf{(b)} Sub-09 (moderate protanomaly, Machado fit): the arc
  compresses from $360^\circ$ to $\sim 96^\circ$; c4 (green), c5
  (cyan), c6 (blue) merge at $\sim 282^\circ$ (red oval) and
  separation-based fallback is required (red crossing arrows).
  Bottom table: only the cortical 2-component model preserves
  bijectivity at both severity levels, making it the preferred filter
  substrate.}
  \label{fig:feasibility}
\end{figure*}

\subsection{Specificity calibration and cross-criterion validation}
\label{sec:results:specificity}

An individual-case inference pipeline with 1--2 free parameters over
8 data points could plausibly achieve nominal significance for any
participant. We therefore quantified this risk empirically by running
the same pipeline on all 7 HC participants, treating each as if they
were a CVD candidate.

\paragraph{HC false-positive rates exceed nominal $\alpha$.}
\cref{fig:hcfpr,tab:fpr} report per-model FPRs under the same
permutation test used for CVD. At $\alpha = 0.05$, 15 of 21
model--subject tests reach significance (Machado 3/7 $= 43\%$; R+C
5/7 $= 71\%$; 2-component 7/7 $= 100\%$), compared with an expected
$1.05/21$ under the global null. A diagnostic analysis identifies two
compounding drivers (\cref{app:supporting}): (i) per-HC baseline
$\rho$ and fit improvement $\Delta\rho$ correlate at $r = -0.894$
across HCs, i.e.\ HC participants with weaker baseline interpolation
have disproportionately more room for spurious gain; and
(ii) rank-based loss on 8 points admits a permutation p-value that
decays rapidly with even mild monotone structure, and the 2-component
model's extra 2 DOF convert this into near-guaranteed rejection of
the null for any individual.

This is a genuine methodological limitation and we do \emph{not} claim
population-level specificity for the pipeline. What we claim is
individual-case specificity against three anchors that the HC FPR
analysis does \emph{not} invalidate:

\begin{enumerate}
\item \textbf{Within-study normal control.} Sub-10 is
non-significant under every cone-level test (Machado/R+C
$p = 0.559$; 2-component $p = 0.058$ marginal). This is the correct
specificity comparator for each CVD case: the HC FPR analysis asks
"can the pipeline falsely flag healthy participants?" and answers
\emph{yes}, but the within-study deutan-axis control asks "does it
flag the participant for whom we have \emph{a priori} reason to expect
no effect?" and answers \emph{no}.
\item \textbf{Physiological interpretability of parameters.} Sub-09's
fitted $\Delta\lambda = 13.5$~nm and Sub-08's
$\Delta\lambda = 2.0$~nm both sit within the 2--12~nm anomalous-cone
range reported by \citet{neitz2011} and compatible with the Ishihara
severity stratification in \cref{sec:results:behavior}. HC fits, by
contrast, distribute uniformly across the grid with no physiological
anchor (\cref{app:supporting}).
\item \textbf{Cross-criterion convergence.} The 2-component model is
the only model that reaches significance for each CVD participant
under both the per-hue LOCO criterion and the pairwise V1 $\Delta$RDM
criterion (Sub-08 LOCO $p = 0.004$, $\Delta$RDM CI excludes 0;
Sub-09 LOCO $p = 0.035$, $\Delta$RDM $p = 0.007$;
\cref{tab:full}). No HC participant achieves cross-criterion
significance in both metrics.
Moreover, the fitted $\beta_s$ from V1 $\Delta$RDM converges at
$\sim 21.5^\circ$ across both CVD cases, matching the $21.4^\circ$
independently reported by \citet{emery2021} from anomalous-trichromat
hue scaling. This quantitative cross-cohort match is not available
from within-sample grid flexibility.
\end{enumerate}

We take (1)--(3) together to establish individual-case specificity for
the two CVD participants at hV4 LOCO, while being transparent that an
uncalibrated screening interpretation of this pipeline on healthy
populations would be unreliable.

\begin{figure*}[tb]
  \centering
  \includegraphics[width=0.95\textwidth]{Figures/generated/fig_hc_specificity.pdf}
  \caption{\textbf{HC specificity calibration.}
  \textbf{(a)} Per-HC $p$-values for each of the three distortion
  models (log scale). Red dashed line: $\alpha = 0.05$; shaded zone:
  nominal false-positive region. Observed 15/21 tests inside the FP
  zone vs.\ expected 1.05/21 under the global null.
  \textbf{(b)} Aggregate FPR by model. Nominal $\alpha = 5\%$
  (gray dashed) and the $1/7$ single-subject chance level
  (gray dotted) are shown for reference. Model flexibility (DOF) and
  rank-based loss jointly inflate the FPR; this is a property of the
  per-case fitting procedure, not of the CVD inference
  (\cref{sec:results:specificity}).}
  \label{fig:hcfpr}
\end{figure*}

\begin{table}[tb]
  \caption{HC false-positive rates under the same grid + permutation
  fitting procedure used for CVD (hV4 LOCO, $\alpha = 0.05$). This
  calibration is used to contextualise CVD significance rather than to
  claim population-level specificity (see
  \cref{sec:results:specificity}).}
  \label{tab:fpr}
  \centering
  \small
  \begin{tabular}{lccc}
    \toprule
    Model & DOF & HC significant & FPR \\
    \midrule
    Machado     & 1 & 3/7 & 43\% \\
    R+C         & 2 & 5/7 & 71\% \\
    2-Component & 2 & 7/7 & 100\% \\
    \midrule
    Total       &   & 15/21 & 71\% \\
    \bottomrule
  \end{tabular}
\end{table}

\paragraph{Cross-ROI replication of the 2-component fit.}
Although hV4 is the primary ROI, the 2-component model replicates in
earlier areas. Sub-08 V1 LOCO reaches $p = 0.001$ ($\beta_s = 50^\circ$,
$\beta_c = -14^\circ$) --- the strongest LOCO effect in the pipeline.
Sub-09 V1 $\Delta$RDM reaches $p = 0.007$ in the independent pairwise
geometry criterion. Sub-10 fails both V1 tests
($p_{\text{LOCO}} = 0.058$ marginal, $p_{\Delta\text{RDM}} = 0.884$).
The full cross-ROI $\times$ cross-criterion matrix is provided in
\cref{tab:full}.

% --- Behavioral validation subsection placeholder (Phase 2 ongoing) ---
% This subsection will be completed once the 2AFC pre-image-filter
% experiment data are acquired (Option B, within 1 month).
%
% \subsection{Behavioural validation of the derived filter}
% \label{sec:results:behavioralfilter}
% Placeholder — to be filled with same-day JND 2AFC results comparing
% (a) uncorrected display, (b) R+C pre-image filter, (c) 2-component
% pre-image filter. Preregistered prediction: along confusion-axis
% hue pairs (Sub-08: c2/c3–c7; Sub-09: c5/c6 merged cluster), only
% the cortical 2-component filter will significantly reduce JND
% relative to uncorrected; the retinal R+C filter will either fail to
% improve (Sub-09, arc compression) or improve differentially
% (Sub-08).

%================================================================
% Figure/table cross-references used:
%   fig:main          — Figure 1 (pipeline + b/c/d LOCO profiles)
%   fig:feasibility   — Figure 2 (correction feasibility / arc collapse)
%   fig:hcfpr         — Figure 3 / S2 (HC FPR forest + bar)
%   tab:main          — Table 1 (inferred parameters hV4 LOCO)
%   tab:fpr           — Table 2 (HC FPR calibration)
%   tab:full          — Supplementary table (cross-ROI, cross-criterion)
%   tab:loro          — Supplementary table (LORO preserved)
%================================================================
